% libre-compendium
%
% Copyright (C) 2026 Richard Qin
%
% This document is dual-licensed under either:
% - GNU Free Documentation License, Version 1.3 or any later version
%   published by the Free Software Foundation; with no Invariant Sections,
%   no Front-Cover Texts, and no Back-Cover Texts (see COPYING.DOC); or
% - Creative Commons Attribution-ShareAlike 4.0 International License
%   (see LICENSE-CC-BY-SA-4.0).

% Exercises migrated from problems06.tex and solutions06.tex.
\begin{exercise}[level=basic]
证明 Logistic 函数满足
\[
\sigma'(z)=\sigma(z)(1-\sigma(z)).
\]
\begin{solution}
\[
\sigma'(z)
=-\frac{1}{(1+e^{-z})^2}(-e^{-z})
=\frac{e^{-z}}{(1+e^{-z})^2}.
\]
又
\[
\sigma(z)(1-\sigma(z))
=\frac1{1+e^{-z}}\frac{e^{-z}}{1+e^{-z}},
\]
两式相等。
\end{solution}
\end{exercise}

\begin{exercise}[level=basic]
对二分类交叉熵，证明
\[
\frac{\partial\ell}{\partial z}=\sigma(z)-y.
\]
\begin{solution}
对
$\ell=-y\log\sigma(z)-(1-y)\log(1-\sigma(z))$，
\begin{align*}
\frac{\partial\ell}{\partial z}
&=-y\frac{\sigma'}{\sigma}
+(1-y)\frac{\sigma'}{1-\sigma}\\
&=-y(1-\sigma)+(1-y)\sigma
=\sigma-y.
\end{align*}
\end{solution}
\end{exercise}

\begin{exercise}[level=basic]
验证 Softmax 对平移不变：
\[
\operatorname{softmax}(z+c\mathbf1)
=\operatorname{softmax}(z).
\]
说明这一性质如何用于避免指数溢出。
\begin{solution}
\[
\frac{e^{z_k+c}}{\sum_je^{z_j+c}}
=\frac{e^ce^{z_k}}{e^c\sum_je^{z_j}}
=\frac{e^{z_k}}{\sum_je^{z_j}}.
\]
因此可取 $c=-\max_jz_j$，使所有新的指数自变量不大于零，避免 $e^{z_j}$ 上溢。
\end{solution}
\end{exercise}

\begin{exercise}[level=basic]
对 $L(\theta)=\theta^2$，写出固定学习率 $\eta$ 的梯度下降迭代，并确定使 $\theta_t\to0$ 的 $\eta$ 取值范围。
\begin{solution}
$L'(\theta)=2\theta$，所以
\[
\theta_{t+1}=(1-2\eta)\theta_t,
\qquad
\theta_t=(1-2\eta)^t\theta_0.
\]
对任意初值收敛到零的充要条件是
\[
|1-2\eta|<1,
\]
即 $0<\eta<1$。
\end{solution}
\end{exercise}

\begin{exercise}[level=standard]
设 $z=Wx+b$、$\widehat y=\operatorname{softmax}(z)$，损失为交叉熵。求 $\nabla_W\ell$、$\nabla_b\ell$ 和 $\nabla_x\ell$。
\begin{solution}
记 $\delta=\widehat y-y$。由 Softmax 与交叉熵的合成梯度，
\[
\nabla_z\ell=\delta.
\]
由 $dz=(dW)x+Wdx+db$，得
\[
\nabla_W\ell=\delta x^{\mathsf T},
\qquad
\nabla_b\ell=\delta,
\qquad
\nabla_x\ell=W^{\mathsf T}\delta.
\]
\end{solution}
\end{exercise}

\begin{exercise}[level=standard]
证明若神经网络各层都使用恒等激活，则任意多层网络都等价于单个仿射变换。
\begin{solution}
对两层有
\[
W_2(W_1x+b_1)+b_2
=(W_2W_1)x+(W_2b_1+b_2),
\]
仍是仿射变换。假设前 $k$ 层的复合为 $Ax+c$，再复合一层 $Wx+b$ 得
$WAx+(Wc+b)$，仍是仿射变换。归纳完成证明。
\end{solution}
\end{exercise}

\begin{exercise}[level=standard]
对两层网络
\[
a=\phi(W_1x+b_1),
\qquad
\widehat y=W_2a+b_2,
\qquad
\mathcal L=\frac12\lVert\widehat y-y\rVert^2,
\]
写出 $W_1,b_1,W_2,b_2$ 的梯度。
\begin{solution}
记
\[
\delta_2=\widehat y-y,
\qquad
\delta_1=(W_2^{\mathsf T}\delta_2)\odot\phi'(W_1x+b_1).
\]
则
\[
\nabla_{W_2}\mathcal L=\delta_2a^{\mathsf T},
\qquad
\nabla_{b_2}\mathcal L=\delta_2,
\]
\[
\nabla_{W_1}\mathcal L=\delta_1x^{\mathsf T},
\qquad
\nabla_{b_1}\mathcal L=\delta_1.
\]
\end{solution}
\end{exercise}

\begin{exercise}[level=standard]
设数据下标 $i$ 在 $\{1,\ldots,N\}$ 中均匀抽取。证明单样本梯度 $g=\nabla\ell_i(\theta)$ 是完整经验风险梯度的无偏估计。
\begin{solution}
由均匀抽样，
\[
\mathbb E[g]
=\sum_{i=1}^N\frac1N\nabla\ell_i(\theta)
=\nabla\left(\frac1N\sum_{i=1}^N\ell_i(\theta)\right)
=\nabla L(\theta).
\]
因此单样本梯度是无偏估计。
\end{solution}
\end{exercise}

\begin{exercise}[level=advanced]
设 $\widehat y=\operatorname{softmax}(z)$。证明
\[
J=\diag(\widehat y)-\widehat y\widehat y^{\mathsf T}
\]
是半正定矩阵，并证明 $J\mathbf1=0$。
\begin{solution}
对任意 $v\in\R^K$，
\begin{align*}
v^{\mathsf T}Jv
&=\sum_i\widehat y_iv_i^2
-\left(\sum_i\widehat y_iv_i\right)^2.
\end{align*}
这是以 $\widehat y_i$ 为概率权重的随机变量 $v_i$ 的方差，因而非负，所以 $J$ 半正定。又
\[
J\mathbf1=\widehat y-widehat y(\widehat y^{\mathsf T}\mathbf1)
=\widehat y-widehat y=0.
\]
这与 Softmax 对整体平移不变一致。
\end{solution}
\end{exercise}

\begin{exercise}[level=advanced]
从反向传播的连乘结构说明：若每层权重矩阵谱范数不超过 $\rho$，激活导数绝对值不超过 $\gamma$，则跨越 $k$ 层传播的梯度范数至多被放大 $(\rho\gamma)^k$ 倍。
\begin{solution}
每向前跨过一层反向传播时，梯度先乘以权重矩阵转置，再与激活导数逐元素相乘。对角导数矩阵的谱范数不超过 $\gamma$，权重矩阵转置的谱范数不超过 $\rho$。由谱范数的次乘性，每层至多乘以 $\rho\gamma$，跨越 $k$ 层后至多乘以
$(\rho\gamma)^k$。若 $\rho\gamma<1$，该上界指数衰减；若它显著大于 $1$，则存在梯度爆炸风险。
\end{solution}
\end{exercise}

